%% kapa: exam=1 round=35 year=2024 period=2 form=A subject=appraisal_law subject_ko=감정평가관계법규 seq=20 number=20
%% source: Q-Net 감정평가사 기출문제 (cst00309) artlSeq=5213229 file=raw/1008/5213229/2교시 문제지 원본.pdf
%% conversion: kordoc 4.24.51 -> tools/gen1.py (auto); review=proofread-sonnet; flags=none
\begin{question}{20}
국유재산법상 용어의 정의이다. ( )에 들어갈 내용으로 옳은 것은?
\begin{framedtext}
○ ( ㄱ )(이)란 국가 외의 자가 제5조제1항 각 호에 해당하는 재산의 소유권을 무상으로 국가에 이전하여 국가가 이를 취득하는 것을 말한다. \par
○ ( ㄴ )이란 사용허가나 대부계약 없이 국유재산을 사용ㆍ수익하거나 점유한 자에게 부과하는 금액을 말한다. \par
○ 총괄청이란 ( ㄷ )을 말한다.
\end{framedtext}
\begin{choices}
  \item ㄱ: 기부채납, ㄴ: 부담금, ㄷ: 중앙관서의 장
  \item ㄱ: 무상양도, ㄴ: 변상금, ㄷ: 기획재정부장관
  \item ㄱ: 기부채납, ㄴ: 변상금, ㄷ: 기획재정부장관
  \item ㄱ: 무상양도, ㄴ: 변상금, ㄷ: 중앙관서의 장
  \item ㄱ: 기부채납, ㄴ: 부담금, ㄷ: 기획재정부장관
\end{choices}
\end{question}
\answer{3}
